% --- Back cover ---
% Large PERSONIX wordmark, no logo image. Cream paper inherited globally.
% Absolute (tikz overlay) positioning so the footer sticks to the bottom.
\clearpage
\thispagestyle{empty}
\begin{tikzpicture}[remember picture, overlay]
  % --- Wordmark ---
  \node[anchor=north, font=\sffamily\bfseries\fontsize{48}{52}\selectfont]
    at ([yshift=-26mm]current page.north) {PERSONIX};
  \node[anchor=north, font=\rmfamily\itshape\large]
    at ([yshift=-44mm]current page.north) {Thay đổi không khoan nhượng};
  \node[anchor=north, font=\rmfamily\small,
        text width=\paperwidth-50mm, align=center]
    at ([yshift=-54mm]current page.north)
    {Một mạng lưới danh tiếng phi tập trung, không thể kiểm duyệt và không thể mua chuộc};
  % --- Body ---
  \node[anchor=north, text width=\paperwidth-50mm, align=justify, font=\small]
    at ([yshift=-72mm]current page.north)
    {%
      Nhà nước vận hành tốt khi giao tiếp còn chậm, các quyết định mang tính
      địa phương, và quyền lực sống trong cùng thị trấn mà nó cai trị. Các mạng
      lưới ngày nay đảo ngược cả ba giả định ấy --- và những thiết chế dựng trên
      chúng không còn phù hợp với thế giới mà chúng cai quản. Cuốn sách này mô tả
      một công cụ thì phù hợp: một mạng lưới danh tiếng phi tập trung, nơi danh
      tính là của bạn để giữ, sự thật có thể kiểm toán công khai, và hối lộ là
      hành vi tự sát về danh tiếng.

      \smallskip
      Không phải một tuyên ngôn. Không phải một dự báo. Một bản thiết kế khả dụng
      cho việc kẻ kế tục nhà nước có thể được dựng nên thế nào --- một cách tự
      nguyện, từng lời tuyên bố một.
    };
  % --- Footer (anchored to the bottom of the page) ---
  \node[anchor=south, font=\sffamily\footnotesize,
        text width=\paperwidth-50mm, align=center]
    at ([yshift=12mm]current page.south)
    {%
      Pavel Kudrna\quad\textbullet\quad Praha, 2026\par
      \href{https://personix.org}{personix.org}\par
      Cấp phép theo Creative Commons Attribution-ShareAlike 4.0 International
      (CC BY-SA 4.0).
    };
\end{tikzpicture}%
\mbox{}%
\clearpage
